\chapter{基因究竟是什么}

\begin{kaichang}
找一张你家的全家福照片，再找一张街头塞给你的、或者手机里刷到的“基因检测”宣传单。宣传单上的词你一定眼熟：检测“运动基因”“肥胖基因”“天赋基因”，还有“酒精代谢基因”“脱发基因”。照片里的你，眉眼像爸爸，脸型像妈妈，可有些地方谁都不像。

现在请你猜一个问题，不用查资料，就凭直觉写下来：如果有人说“测一测你的某个基因”，他到底要测什么？是测你细胞里某一种蛋白质的含量？是数一数你有多少条染色体？还是在你那本 30 亿字母的 DNA 长卷里，找到某一小段，读出它的字母？

“基因”可能是现代科学里被用得最多、被说得最含混的词。广告里有它，新闻里有它，长辈说“这孩子随他爸，基因好”时也有它。可要是追问一句：一个基因，作为一个实实在在的东西，究竟是什么？能答上来的人并不多。这一章就来回答这个问题。到章末，我们会回到这张宣传单，看看“检测一个基因”这句话到底是什么意思。

先看几个每天都能碰到的句子：“他有运动天赋，是基因好”“这个毛病是遗传的”“转基因食品能不能吃”“警察通过 DNA 锁定了嫌疑人”。这几句话里的“基因”指的其实是不同层面的东西：一个性状、一种疾病倾向、一段被人工转移的序列、一堆用来比对身份的字母差异。一个词被这么多人拿去干这么多活，含混几乎是注定的。所以本章的任务不是再给你一个要背的定义，而是带你把“基因”拆到能看清零件的程度。
\end{kaichang}

\section{人人都说基因，基因在哪儿}

先说能看见的。你和父母长得像，又不完全像；豌豆花有紫有白，种子有圆有皱（第 63 章）；同一窝小猫，毛色各不相同。这些都是肉眼可见的事实，古人早就知道“种瓜得瓜”。但“像”这件事本身是说不清楚的——它像一团雾，笼罩在每一个家庭上，谁都能感觉到，谁也抓不住。

孟德尔的伟大之处，是把这团雾拆成了可以数的东西。他不问“为什么孩子像父母”这种大而无当的问题，而是挑了豌豆里界限分明的几对性状：花不是紫就是白，种子不是圆就是皱。然后他去数数：杂交后代里，紫花和白花各有多少株？比例是多少（第 63 章）？数着数着，雾散开了，露出了后面有板有眼的数学规律。于是他推测：一定有某种“遗传因子”，一个性状一对，从父母各得一个。

注意，孟德尔从头到尾没见过任何“因子”。他不知道因子是分子还是别的什么，不知道它住在细胞里哪个角落。他只有一个抽象的概念：\textbf{控制一个性状的、可以拆开又组合的、按概率传递的单位}。这是基因的第一个定义——一个纯粹从数数中推断出来的定义。它有点像第 7 章里道尔顿的原子：先假设有这种单位，因为它能解释称量出来的数字，至于它长什么样，留给后来的人。

后来的故事你在前几章已经走过了：染色体被发现，因子被安顿在染色体上（第 64 章）；DNA 被证明是遗传物质（第 65 章）；双螺旋揭示了信息就写在碱基的排列里（第 66 章）；复制、转录、翻译的机器被一一拆开（第 67 到 69 章）；开关这些信息的调控系统也浮出了水面（第 70 章）。现在，链条走到这里，我们终于可以给那个抽象的“因子”安一个物质的地址了。

\section{把镜头推进到 DNA 长链}

把细胞核里的一条染色体想象成一卷超长的磁带，磁带上用 A、T、G、C 四种字母密密麻麻写满了字。人的体细胞里有 46 卷这样的磁带（23 对），加起来约 30 亿个字母，合称\textbf{基因组}——你的全部遗传信息的总和。这套长卷你有两份，一份来自父亲，一份来自母亲，所以绝大多数位置你都有两个抄本。

那么一个基因在哪里？就在长卷的\textbf{某一个固定位置}上。这个位置有个专门的名字，叫\textbf{基因座}（“座”就是座位的座）。比如，控制你是否能消化乳糖的乳糖酶基因，就坐在第 2 号染色体上一个固定的座位上。每个人的这个座位上都有一段 DNA，但字母可以略有不同：你的是这个版本，隔壁同学的可能是另一个版本。同一基因座上的不同版本，叫\textbf{等位基因}。第 63 章里孟德尔的“显性、隐性因子”，说的其实就是同一座位上的不同等位基因——显性只是杂合时的表现关系，绝不意味着它更强、更常见或者更好。

把这几个词的关系钉牢：

\begin{itemize}[itemsep=2pt]
\item \textbf{基因组}：全套长卷，约 30 亿字母，两套抄本。
\item \textbf{基因座}：长卷上的一个固定座位、固定地址。
\item \textbf{基因}：坐在那个座位上、能产生某种功能产物的一段序列。
\item \textbf{等位基因}：同一座位上这段序列的不同版本。
\end{itemize}

注意第三个定义，它跟孟德尔的定义已经不一样了。孟德尔说基因是“控制一个性状的因子”，而今天我们说基因是“\textbf{一段能产生功能产物的 DNA 序列}”。产物通常是蛋白质，也可以是直接干活的 RNA。为什么要改定义？因为科学往下挖了一层之后发现，“控制性状”这个说法太间接、太容易误导了——从一段序列到一个性状，中间隔着转录、翻译、蛋白质网络、细胞和整个身体（第 73 章会专门走这段路）。而“产生功能产物”是可以直接在实验室里检验的：这段序列转录出了什么？翻译出了什么？这个产物在细胞里干什么？

再体会一下“等位基因”有多日常。你的血型就是一个基因座上的版本问题：ABO 血型由第 9 号染色体上的一个基因座决定，那里流传着一个能造糖基转移酶的基因，酶的不同版本会把不同的糖链装到红细胞表面。A 版本装上 A 型糖，B 版本装上 B 型糖，而 O 版本是一次“读不下去”的改变，酶失去了功能，什么糖都不装。你有两个抄本，所以可能是 AA、AO、BB、BO、AB 或 OO——六种组合，四种血型。所谓“显性”，在这里只是“装得上糖就压过装不上”，A 和 B 之间则是平起平坐的共显性（第 63 章）。没有谁更“高级”，只有版本不同。

\begin{qiaoliang}
基因组这本“书”到底有多大？做个数量级估算。人的单套基因组约 $3\times 10^{9}$ 个字母。假如用排版这本书的字号把它们印出来，一页大约能排 3000 个字母，那么：
\[3\times 10^{9} \div 3000 = 10^{6}\ \text{页}\]
一百万页。按一本书 500 页算，相当于 2000 本书，塞满一个小型图书馆——而这全部装在直径只有约 \ud{6}{\mu m} 的细胞核里（第 66 章讲过它怎样被压缩打包）。
再算基因的占比。人类基因组里，编码蛋白质的基因大约 2 万个，听起来很多，可是真正翻译蛋白质的编码序列加起来只占总长度的 1.5\% 左右。也就是说，那 2000 本书里，只有 30 本是“正文”，其余是调控开关、转录成 RNA 直接工作的序列、旧零件、重复段落，以及至今功能不明的部分。这个比例本身就是一条重要信息：基因组远不是一本“基因排排坐”的字典。
\end{qiaoliang}

\section{一个基因的内部长什么样}

现在把一个基因单独拿出来，放大看它的内部结构。以真核生物（包括你）的典型基因为例，它至少包含两类区域。

第一类是\textbf{调控区}，在基因的“车头”附近，最出名的是启动子（第 68 章）。它不编码任何产物，却是 RNA 聚合酶和各种转录因子的停靠台，决定这个基因在什么细胞、什么时候、以多大的音量被读取。第 70 章讲过，基因不是一直开着的灯，而开关和调光旋钮就装在这里。调控区出了错，产物的氨基酸序列可以完全正常，却因为“在错误的时间、错误的细胞里被制造”而惹出疾病——单看编码序列是看不出这种问题的。

第二类是\textbf{转录区}，是真正被抄写成 RNA 的部分。其中最终被翻译成蛋白质的那段叫\textbf{编码区}——密码子在这里被三个三个地读成氨基酸（第 69 章）。而在真核生物里，转录区内部还分两种段落：\textbf{外显子}是最终会拼进成品 mRNA 的段落，\textbf{内含子}是抄下来之后又被剪掉、不进入成品的段落（剪接过程见第 68 章）。编码区由外显子拼接而成，中间可能隔着许多内含子；同一个基因通过不同的剪接方式，可以产出好几种不同的蛋白质——这也是为什么“一个基因对应一种蛋白质”只是历史上的简化。

用一张示意图把这些零件摆在一起（这是人为的表示法，方块的宽窄不代表真实长度，真实的内含子往往比外显子长得多）：

\begin{center}
\begin{tikzpicture}[x=1cm,y=1cm]
\draw[thick] (0,0) -- (12,0);
\draw[fill=gray!30] (0.8,-0.3) rectangle (2.0,0.3);
\draw[fill=black!60] (3.0,-0.3) rectangle (4.2,0.3);
\draw[fill=gray!15] (4.2,-0.15) rectangle (6.2,0.15);
\draw[fill=black!60] (6.2,-0.3) rectangle (7.6,0.3);
\draw[fill=gray!15] (7.6,-0.15) rectangle (9.2,0.15);
\draw[fill=black!60] (9.2,-0.3) rectangle (10.4,0.3);
\node at (1.4,0.65) {\small 启动子};
\node at (3.6,0.65) {\small 外显子};
\node at (5.2,0.6) {\small 内含子};
\node at (6.9,0.65) {\small 外显子};
\node at (8.4,0.6) {\small 内含子};
\node at (9.8,0.65) {\small 外显子};
\draw[->,thick] (0.1,0) -- (0.8,0);
\node at (0.35,-0.5) {\small 转录起点};
\end{tikzpicture}
\end{center}

所以，“一个基因”不是长卷上一段边界分明的、均质的文字，而是一个\textbf{带开关、带剪辑标记的功能区域}。一段 DNA 算不算基因，不在于它有多长，而在于细胞会不会拿它做出有用的产物，以及它有没有配套的读取开关。

\section{这个定义是被一场实验磨出来的}

“基因是产生功能产物的序列区域”这句话听起来平淡，却是被一代代实验反复打磨出来的。其中最漂亮的一仗，是“一个基因一种酶”假说。

\begin{zhengju}
1940 年代，美国科学家比德尔和塔特姆选了粉红色面包霉菌（粗糙脉孢菌）做研究对象。这种霉菌有个好处：野生型只要吃简单的糖、盐和一种维生素就能活，因为它自己会制造各种氨基酸。他们的思路很直接：如果基因真的是“干活的指令”，那么弄坏一个基因，就应该弄坏一种具体的制造能力。

他们用 X 射线照射霉菌孢子，人为制造突变（第 72 章会细讲突变的种类），然后挨个检查后代：这株还能不能在“极简食谱”上活？绝大多数突变株照样能活，说明被打坏的地方无关紧要。但总有几株突然“挑食”了——必须在食谱里额外添加某一种物质，比如精氨酸，才能生长。更妙的是，他们接着做精细排查：给这株加精氨酸能活，加瓜氨酸行不行？加鸟氨酸行不行？结果有的突变株加鸟氨酸就能活，有的必须加瓜氨酸，有的非精氨酸不可。这说明制造精氨酸是一条多步流水线，不同的突变株卡在流水线的不同工位上——\textbf{每个工位对应一种酶，每种酶对应一个被打坏的基因}。1941 年，他们据此提出：一个基因的职责，大体是决定一种酶（后来推广为一种蛋白质）。两人因此获得 1958 年诺贝尔生理学或医学奖。

这个故事里也有后来人修掉的毛边。“一个基因一种酶”很快被证明是简化：有的蛋白质由好几个基因的产物拼装而成；有的基因产物是 RNA 而不是蛋白质；可变剪接让一个基因能造出多种蛋白质（第 68 章）。科学定义就是这样：先立一个能打仗的版本，再在证据里一点点磨出更准确的版本。本章开头你看到的那个定义——“产生功能产物的序列区域”——就是磨到今天的结果。
\end{zhengju}

\section{定义边界上的三个“麻烦户”}

磨到今天的定义依然有不服帖的地方。这不是科学的失败，而是大自然本来就不按人类的名词表生长。认识三个“麻烦户”，你对基因的理解才算完整。

\textbf{非编码 RNA 基因}。有些序列转录出来的 RNA 不去翻译蛋白质，而是直接干活：tRNA 当翻译的适配器，rRNA 是核糖体的骨架和催化核心（第 69 章），还有各种调控 RNA 参与开关别的基因（第 70 章）。按“控制性状的因子”这个老定义，它们很难算基因；按“产生功能产物的序列区域”这个新定义，它们理直气壮地是基因。定义的改变，正是因为发现了它们。

\textbf{重叠基因}。在一些病毒里，同一段 DNA 序列可以被按不同的读法、不同的起点读取，编码两种甚至三种不同的蛋白质——就像一句话换一种断句方式就变了意思。这时“一个基因”的边界在哪里？序列重叠在一起，你没法在磁带上画出一条不打架的界线。

\textbf{假基因}。基因组里还散落着许多“旧零件”：它们长得明显像某个真基因，却因为积累了破坏性的改变而不再生产产物。它们是进化留下的废弃文件，有时会被重新启用或拆解利用（第 79 章会讲基因组怎样重复、借用和改造）。算不算基因？得看你强调“有产物”还是“有来历”。

所以今天的生物学家用“基因”这个词时，心里其实带着一把软尺子：\textbf{能产生功能产物、在进化中被维持的那段序列区域}。边界模糊的个案交给上下文去处理。名词是为人服务的，不是自然必须遵守的。

\begin{wuqu}
“人的基因组就是一本人体说明书，每个基因对应书里的一个零件条目。”——这个比喻流传极广，但它会把你带沟里去。首先，基因组里没有任何“条目”直接描述性状；基因描述的是产物（蛋白质或 RNA），从产物到性状中间隔着整个生命（第 73 章）。其次，说明书是为组装一台机器写的，而基因组没有人“写”，它是几十亿年试错攒下的、充满重复段落和废弃文件的档案库。还有最实际的一条：同一份基因组在不同细胞里被读出的内容完全不同——神经细胞和肌肉细胞用着同一套“书”，却读出了两种人生（第 70 章）。把基因组想成一座堆满配方的老厨房更贴切：配方本身不会做菜，做出什么菜，取决于谁掌勺、用哪张配方、火候如何。
\end{wuqu}

\section{基因在网络里，不在真空里}

本章给基因划了地址、看了内部、修了定义，但必须立刻补上一句防止误解的话：\textbf{单个基因几乎从来不单独决定什么}。

一个基因的产物进入细胞后，立刻卷入网络：乳糖酶要在肠道细胞里、在正确的时间被制造出来才有用；它制造多少，受调控区控制（第 70 章）；它工作好坏，取决于有没有被正确折叠、运输（第 69 章）。而一个看得见的性状——身高、肤色、血型、能不能喝牛奶——往往是许多基因的产物层层协作、再叠加环境和经历之后的结果。把“一个基因”从网络里剪出来问“它决定什么”，就像从交响乐团里拽出一位乐手问“这首曲子是不是你演奏的”。

这也解释了为什么“等位基因”才是遗传学里真正流通的货币。谈论“乳糖酶基因”时，大多数人都有它；真正造成人与人差别的，是那个座位上的\textbf{版本}不同：有人一辈子能消化乳糖，有人成年后肠道里这个基因被调低了音量（第 70 章、第 74 章）。基因检测报告里写的“你携带某某型”，翻译过来就是：在某个基因座上，你的两个抄本是哪两个版本。版本从哪来？从突变来——那是下一章的主角。

还要说清一件经常被弄反的事：环境和网络不是基因之外的“干扰项”，而是基因工作方式的一部分。同一种植物，种在背阴处和向阳处，同一套基因会被读出不同的生长策略；同一种蜜蜂的幼虫，吃普通花蜜的发育成工蜂，吃蜂王浆的发育成蜂王——差别不在序列，而在序列被怎样读取（第 74 章的表观遗传会细讲其中一种机制）。基因为可能性划定范围，环境和调控在这个范围内做出选择。把基因想成骰子的点数上限，把环境和发育想成那只掷骰子的手，比“基因决定一切”要准确得多。

\begin{shiyan}
动手画一画你自己的“版本地图”。准备一张纸、一支笔，再做一点家庭作业：
\begin{enumerate}[itemsep=1pt]
\item 列出三个人：你自己、父亲、母亲（或两位与你朝夕相处的长辈）。
\item 选三个容易观察的性状版本：能否卷舌、有无酒窝、拇指末节能否明显后仰。对每个人逐项记录“有/无”或“能/不能”。
\item 把结果画成一张三行三列的表，在每个格子里标注你观察到的版本。
\item 观察点：哪些版本你和父母一方一致？有没有你不像任何一方的？想一想，你的两个抄本分别来自谁，表里的“有/无”背后可能藏着哪几种等位基因组合（别忘了隐性版本可以不露面，第 63 章）。
\end{enumerate}
安全提示：这是纯观察活动。注意不要把“能不能卷舌”当成任何能力或健康的指标——它只是一个无伤大雅的版本差异；也不要仅凭这些特征去做任何亲子关系的推断，单一性状版本说明不了那类问题。
\end{shiyan}

\section{回到那张宣传单}

现在可以拆开章头那张宣传单了。“检测你的酒精代谢基因”，实际做的事是：取一点你的口腔细胞，提取 DNA，在乙醛脱氢酶基因的那个基因座上，读出你两个抄本的字母，看你携带的是哪个等位基因版本。如果你的版本让这个酶的活性偏低，乙醛分解得慢，喝酒容易脸红——报告预测的是\textbf{一种化学反应速率上的倾向}，不是“你能喝多少酒”的判决书，因为实际酒量还受体重、性别、饮酒习惯等众多因素影响。

“天赋基因”“运动基因”就更值得警惕了。只要这个性状是许多基因加环境共同塑造的（几乎一定如此，第 73 章会拆开讲），任何“测一个基因定天赋”的说法都超出了证据。下次再看到这类广告，你可以问三个问题：测的是哪个基因座？这个版本影响的是哪个产物的什么功能？从产物到广告说的那个“天赋”，中间的因果链拿出证据了吗？

顺带再看一眼那张全家福。现在你明白了：你从父亲和母亲那里各拿到一套长卷，绝大多数座位上坐着相同的基因，但相当一部分座位上，两个抄本是不同的版本。你像爸爸的地方，可能恰好是某个座位上显性版本露了面；你谁都不像的地方，可能是两个隐性版本碰巧配了对（第 63 章），也可能是某个多基因性状落在了一个父母都没到过的组合上。相似和不同，用的其实是同一套规则。

\section{伏笔：版本从哪里来}

这一章我们反复说“同一座位上的不同版本”，却一直回避了一个问题：这些版本最初是怎么冒出来的？为什么人类这个座位上普遍是这个字母，偶尔会出现另一个字母？字母变了之后，产物会跟着变吗——是变坏、变好，还是根本无所谓？

答案藏在一个既不神秘也不浪漫的过程里：抄写会出错，射线会打断，病毒会插足——而细胞有一整套修理工，修得极好，却不是万无一失。第 72 章，我们去看“突变”这位被误解最深的角色：它既是绝大多数遗传病的来源，也是进化手里唯一的原材料。

\begin{tiaozhan}
（跳过不影响主线）把基因座想成一条数轴上的区间，你会遇到一个微妙的数学问题：一个基因座的两个抄本，各自是 30 亿字母里的一段。如果每个字母位置平均约每 $10^{8}$ 份配子 DNA 中才出一次复制错误（第 67 章的校对精度），那么一个长度为 $3\times 10^{4}$ 个字母的基因，在一代之内发生一次新错误的概率大约是 $3\times 10^{4}\times 10^{-8}=3\times 10^{-4}$，即约万分之三。听起来很小，但人体有约 2 万个基因，加起来，平均每个人出生时都带着几十处父母都没有的全新字母差异。群体层面再乘以人口数量，几乎所有可错的字母，每一代都有人错了一遍。这就是为什么“有没有新版本”从来不是概率问题，而是时间问题——第 72 章和第 77 章会顺着这个计算继续往下走。
\end{tiaozhan}

\section*{练一练}

\begin{sikaoti}
\item 用你自己的话区分基因组、基因座、基因、等位基因这四个词，并各举一个本章之外的例子（可以编：比如假设某植物的花色基因座）。
\item 把“基因是产生功能产物的序列区域”和孟德尔时代“基因是控制性状的因子”这两个定义并排写，说说新定义赢在哪里、又为什么仍需要“网络和环境”来补充。
\item 画一幅示意图：一条染色体长卷上标出两个基因座，在其中一个基因座上画出两个人的两个抄本，用不同底纹表示不同的等位基因版本。图旁注明：这是表示法，真实 DNA 没有“底纹”。
\item 比德尔和塔特姆的实验里，一株突变霉菌“加精氨酸能活、加瓜氨酸也能活、加鸟氨酸不能活”。根据流水线的顺序推理：这株霉菌的哪一个工位坏了？（提示：能被补上的物质在坏工位的下游。）
\item 为什么“重叠基因”和“非编码 RNA 基因”会让“一个基因是什么”变得难定义？这两个麻烦各自挑战了定义里的哪一条？
\item 一张宣传单声称“检测决定数学天赋的基因”。参照本章末尾的三个问题，写一段话向家人解释为什么这种说法超出了证据。
\end{sikaoti}
